\documentclass[12pt]{standalone}
\usepackage{fontspec}
\setsansfont{TeX Gyre Heros}
\newfontfamily\mfallback{DejaVu Sans}
\usepackage{tikz}
\usetikzlibrary{arrows.meta}
% =====================================================================
% 统一符号层（2026-10-20）
% 用法：在 fig_<号>.tex 的 preamble 里 % =====================================================================
% 统一符号层（2026-10-20）
% 用法：在 fig_<号>.tex 的 preamble 里 % =====================================================================
% 统一符号层（2026-10-20）
% 用法：在 fig_<号>.tex 的 preamble 里 \input{../_symbols/symbols.tex}
% 依赖：fontspec（已 \setsansfont{TeX Gyre Heros}）、graphicx 不需要
% 规则见同目录《符号验收清单.md》
% =====================================================================

% ---- 字体角色（全部走 sans = TeX Gyre Heros）----
\newcommand{\SIt}[1]{\textsf{\textit{#1}}}          % 斜体变量
\newcommand{\SBf}[1]{\textsf{\textbf{#1}}}          % 粗体
\newcommand{\SBfIt}[1]{\textsf{\textbf{\textit{#1}}}} % 粗斜体变量
\newcommand{\SNum}[1]{\textsf{\textup{#1}}}         % 直立数字

% ---- 上下标：字号/位置写死，避免数学样式与手调混用 ----
% 审计结论（2026-10-20）：源图下标/底字高比在 0.77–0.92em，0.72em 偏小。
% 默认 0.80em；仍需按源图逐图标定（44.4/44.5 约 0.85–0.92，39/41/46 约 0.77–0.80）。
\newcommand{\FSup}[2]{#1\raisebox{0.30em}{\fontsize{0.72em}{0.82em}\selectfont #2}}
\newcommand{\FSub}[2]{#1\raisebox{-0.30em}{\fontsize{0.80em}{0.92em}\selectfont #2}}

% ---- 分数：数字直立、方头粗杆；杆宽 = max(分子宽,分母宽) + 2pt ----
% 杆厚 4pt（配合 20–26pt 标签）；小字号图可局部调 0.17em
% 2026-10-20 修正：原 \shortstack 负 kern 不随字号缩放（23pt 下总高仅 30pt、杆贴字），
% 改为 offinterlineskip + em 定间距：上间隙 0.12em、下间隙 0.13em
% （37.2 源图 23pt 实测四标签上/下间隙均值 2.6/4.0pt；取近平衡且总高保持 42–43pt）。
\newcommand{\FFrac}[2]{%
  \begingroup
  \setbox0=\hbox{#1}\setbox2=\hbox{#2}%
  \dimen0=\wd0 \ifdim\wd2>\dimen0 \dimen0=\wd2 \fi
  \advance\dimen0 by 2pt
  \leavevmode\vbox{\offinterlineskip
    \hbox to \dimen0{\hss\box0\hss}%
    \kern0.12em
    \hbox{\rule{\dimen0}{4pt}}%
    \kern0.13em
    \hbox to \dimen0{\hss\box2\hss}%
  }%
  \endgroup}

% ---- 真减号（审计：源图减号 15–18px 宽、3–5px 厚；Heros 常规字重偏细）----
% 可选参数 = 宽度（em）：正文 0.58；上标内用 \FMinus[0.72]
\newcommand{\FMinus}[1][0.58]{\rule[0.30em]{#1em}{0.13em}}

% ---- 直撇（审计：44.2 用弯引号、44.3 用直撇，必须统一为直楔撇）----
\newcommand{\FPrime}{\textsf{\textquotesingle}}

% ---- 乘号（审计：46.2 的 × 比源小 25–30%，标签内需按源校准）----
% 用 \FTimes 时在 node 内与底字同字号，并在目检时量宽高。


% ---- 斜杠分数（审计 A：35.1 的 r/3 用 38pt 斜杠手搓）----
\newcommand{\FSlashFrac}[2]{#1\raisebox{-0.18em}{\fontsize{1.35em}{1.35em}\selectfont\textsf{/}}\kern-0.06em #2}

% ---- 交集号 ∩（Heros 无此字，固定回退字号/基线，避免逐图不同）----
\newcommand{\FIntersect}{\raisebox{0.02em}{\fontsize{0.95em}{1em}\selectfont\mfallback{∩}}}

% ---- 括号/开区间（图形层，不是字形）----
% 一律 \draw ... arc[start angle=, end angle=, radius=]：
%   · 端点记号背离它所界的区间；
%   · 左端/右端分别 10× 特写核对；
%   · 闭区间压在轴上时竖脊右移半个线宽；
%   · 同一张图同类记号同一种写法。

% 依赖：fontspec（已 \setsansfont{TeX Gyre Heros}）、graphicx 不需要
% 规则见同目录《符号验收清单.md》
% =====================================================================

% ---- 字体角色（全部走 sans = TeX Gyre Heros）----
\newcommand{\SIt}[1]{\textsf{\textit{#1}}}          % 斜体变量
\newcommand{\SBf}[1]{\textsf{\textbf{#1}}}          % 粗体
\newcommand{\SBfIt}[1]{\textsf{\textbf{\textit{#1}}}} % 粗斜体变量
\newcommand{\SNum}[1]{\textsf{\textup{#1}}}         % 直立数字

% ---- 上下标：字号/位置写死，避免数学样式与手调混用 ----
% 审计结论（2026-10-20）：源图下标/底字高比在 0.77–0.92em，0.72em 偏小。
% 默认 0.80em；仍需按源图逐图标定（44.4/44.5 约 0.85–0.92，39/41/46 约 0.77–0.80）。
\newcommand{\FSup}[2]{#1\raisebox{0.30em}{\fontsize{0.72em}{0.82em}\selectfont #2}}
\newcommand{\FSub}[2]{#1\raisebox{-0.30em}{\fontsize{0.80em}{0.92em}\selectfont #2}}

% ---- 分数：数字直立、方头粗杆；杆宽 = max(分子宽,分母宽) + 2pt ----
% 杆厚 4pt（配合 20–26pt 标签）；小字号图可局部调 0.17em
% 2026-10-20 修正：原 \shortstack 负 kern 不随字号缩放（23pt 下总高仅 30pt、杆贴字），
% 改为 offinterlineskip + em 定间距：上间隙 0.12em、下间隙 0.13em
% （37.2 源图 23pt 实测四标签上/下间隙均值 2.6/4.0pt；取近平衡且总高保持 42–43pt）。
\newcommand{\FFrac}[2]{%
  \begingroup
  \setbox0=\hbox{#1}\setbox2=\hbox{#2}%
  \dimen0=\wd0 \ifdim\wd2>\dimen0 \dimen0=\wd2 \fi
  \advance\dimen0 by 2pt
  \leavevmode\vbox{\offinterlineskip
    \hbox to \dimen0{\hss\box0\hss}%
    \kern0.12em
    \hbox{\rule{\dimen0}{4pt}}%
    \kern0.13em
    \hbox to \dimen0{\hss\box2\hss}%
  }%
  \endgroup}

% ---- 真减号（审计：源图减号 15–18px 宽、3–5px 厚；Heros 常规字重偏细）----
% 可选参数 = 宽度（em）：正文 0.58；上标内用 \FMinus[0.72]
\newcommand{\FMinus}[1][0.58]{\rule[0.30em]{#1em}{0.13em}}

% ---- 直撇（审计：44.2 用弯引号、44.3 用直撇，必须统一为直楔撇）----
\newcommand{\FPrime}{\textsf{\textquotesingle}}

% ---- 乘号（审计：46.2 的 × 比源小 25–30%，标签内需按源校准）----
% 用 \FTimes 时在 node 内与底字同字号，并在目检时量宽高。


% ---- 斜杠分数（审计 A：35.1 的 r/3 用 38pt 斜杠手搓）----
\newcommand{\FSlashFrac}[2]{#1\raisebox{-0.18em}{\fontsize{1.35em}{1.35em}\selectfont\textsf{/}}\kern-0.06em #2}

% ---- 交集号 ∩（Heros 无此字，固定回退字号/基线，避免逐图不同）----
\newcommand{\FIntersect}{\raisebox{0.02em}{\fontsize{0.95em}{1em}\selectfont\mfallback{∩}}}

% ---- 括号/开区间（图形层，不是字形）----
% 一律 \draw ... arc[start angle=, end angle=, radius=]：
%   · 端点记号背离它所界的区间；
%   · 左端/右端分别 10× 特写核对；
%   · 闭区间压在轴上时竖脊右移半个线宽；
%   · 同一张图同类记号同一种写法。

% 依赖：fontspec（已 \setsansfont{TeX Gyre Heros}）、graphicx 不需要
% 规则见同目录《符号验收清单.md》
% =====================================================================

% ---- 字体角色（全部走 sans = TeX Gyre Heros）----
\newcommand{\SIt}[1]{\textsf{\textit{#1}}}          % 斜体变量
\newcommand{\SBf}[1]{\textsf{\textbf{#1}}}          % 粗体
\newcommand{\SBfIt}[1]{\textsf{\textbf{\textit{#1}}}} % 粗斜体变量
\newcommand{\SNum}[1]{\textsf{\textup{#1}}}         % 直立数字

% ---- 上下标：字号/位置写死，避免数学样式与手调混用 ----
% 审计结论（2026-10-20）：源图下标/底字高比在 0.77–0.92em，0.72em 偏小。
% 默认 0.80em；仍需按源图逐图标定（44.4/44.5 约 0.85–0.92，39/41/46 约 0.77–0.80）。
\newcommand{\FSup}[2]{#1\raisebox{0.30em}{\fontsize{0.72em}{0.82em}\selectfont #2}}
\newcommand{\FSub}[2]{#1\raisebox{-0.30em}{\fontsize{0.80em}{0.92em}\selectfont #2}}

% ---- 分数：数字直立、方头粗杆；杆宽 = max(分子宽,分母宽) + 2pt ----
% 杆厚 4pt（配合 20–26pt 标签）；小字号图可局部调 0.17em
% 2026-10-20 修正：原 \shortstack 负 kern 不随字号缩放（23pt 下总高仅 30pt、杆贴字），
% 改为 offinterlineskip + em 定间距：上间隙 0.12em、下间隙 0.13em
% （37.2 源图 23pt 实测四标签上/下间隙均值 2.6/4.0pt；取近平衡且总高保持 42–43pt）。
\newcommand{\FFrac}[2]{%
  \begingroup
  \setbox0=\hbox{#1}\setbox2=\hbox{#2}%
  \dimen0=\wd0 \ifdim\wd2>\dimen0 \dimen0=\wd2 \fi
  \advance\dimen0 by 2pt
  \leavevmode\vbox{\offinterlineskip
    \hbox to \dimen0{\hss\box0\hss}%
    \kern0.12em
    \hbox{\rule{\dimen0}{4pt}}%
    \kern0.13em
    \hbox to \dimen0{\hss\box2\hss}%
  }%
  \endgroup}

% ---- 真减号（审计：源图减号 15–18px 宽、3–5px 厚；Heros 常规字重偏细）----
% 可选参数 = 宽度（em）：正文 0.58；上标内用 \FMinus[0.72]
\newcommand{\FMinus}[1][0.58]{\rule[0.30em]{#1em}{0.13em}}

% ---- 直撇（审计：44.2 用弯引号、44.3 用直撇，必须统一为直楔撇）----
\newcommand{\FPrime}{\textsf{\textquotesingle}}

% ---- 乘号（审计：46.2 的 × 比源小 25–30%，标签内需按源校准）----
% 用 \FTimes 时在 node 内与底字同字号，并在目检时量宽高。


% ---- 斜杠分数（审计 A：35.1 的 r/3 用 38pt 斜杠手搓）----
\newcommand{\FSlashFrac}[2]{#1\raisebox{-0.18em}{\fontsize{1.35em}{1.35em}\selectfont\textsf{/}}\kern-0.06em #2}

% ---- 交集号 ∩（Heros 无此字，固定回退字号/基线，避免逐图不同）----
\newcommand{\FIntersect}{\raisebox{0.02em}{\fontsize{0.95em}{1em}\selectfont\mfallback{∩}}}

% ---- 括号/开区间（图形层，不是字形）----
% 一律 \draw ... arc[start angle=, end angle=, radius=]：
%   · 端点记号背离它所界的区间；
%   · 左端/右端分别 10× 特写核对；
%   · 闭区间压在轴上时竖脊右移半个线宽；
%   · 同一张图同类记号同一种写法。

\begin{document}
\begin{tikzpicture}[x=1pt, y=-1pt, line cap=round, line join=round]

  % ---------- labels (drawn first, white background; upright bold) ----------
  \node[anchor=center,inner sep=0pt,fill=white,font=\fontsize{31}{37}\selectfont\sffamily\bfseries] at (820,206.5)
    {\SBf{B}\FSub{}{1}};
  \node[anchor=center,inner sep=0pt,fill=white,font=\fontsize{31}{37}\selectfont\sffamily\bfseries] at (852,299)
    {\SBf{E}(\SBf{B}\FSub{}{1})};
  \node[anchor=center,inner sep=0pt,fill=white,font=\fontsize{31}{37}\selectfont\sffamily\bfseries] at (816,510.5)
    {\SBf{B}\FSub{}{2}};
  \node[anchor=center,inner sep=0pt,fill=white,font=\fontsize{31}{37}\selectfont\sffamily\bfseries] at (852,612)
    {\SBf{E}(\SBf{B}\FSub{}{2})};

  % ---------- hatch inside B1: backslash family, 13pt spacing ----------
  \begin{scope}
    \clip (390.8,334.5) circle (136.2);
    \foreach \d in {-143,-130,-117,-104,-91,-78,-65,-52,-39,-26,-13,0,13,26,39,52,65,78,91,104,117,130,143}{
      \draw[line width=3.0] ({389.87+\d*0.7071-320},{341.13-\d*0.7071-320}) -- ({389.87+\d*0.7071+320},{341.13-\d*0.7071+320});
    }
  \end{scope}

  % ---------- hatch inside B2: slash family, 13pt spacing ----------
  \begin{scope}
    \clip (258.5,505) circle (120.5);
    \foreach \d in {-130,-117,-104,-91,-78,-65,-52,-39,-26,-13,0,13,26,39,52,65,78,91,104,117,130}{
      \draw[line width=3.0] ({252.97+\d*0.7071-320},{501.97+\d*0.7071+320}) -- ({252.97+\d*0.7071+320},{501.97+\d*0.7071-320});
    }
  \end{scope}

  % ---------- square grid ----------
  \foreach \x in {16,88,164,236,306,378,450,520,590,660,729}{
    \draw[line width=3.4] (\x,16) -- (\x,718);
  }
  \foreach \y in {16,86,155,224,298,368,440,510,580,650,718}{
    \draw[line width=3.4] (16,\y) -- (729,\y);
  }

  % ---------- blue boundaries：每个元素所穿格子的并集边界（6.5pt）----------
  % 左上斜线 (53,118)-(123,187) 所穿的 3 格
  \draw[blue, line width=6.5, line join=miter]
    (16,86) -- (88,86) -- (164,86) -- (164,155) -- (164,224)
    -- (88,224) -- (88,155) -- (16,155) -- cycle;
  % 右下斜线 (606,486)-(656,620) 的 3 格 + E(B2) 格
  \draw[blue, line width=6.5, line join=miter]
    (590,440) -- (660,440) -- (660,580) -- (729,580)
    -- (729,650) -- (590,650) -- cycle;
  % 两圆所穿格子的并集边界（按圆与格子的实际相交计算）
  \draw[blue, line width=6.5, line join=miter]
    (236,298) -- (236,368) -- (164,368) -- (88,368) -- (88,440) -- (88,510)
    -- (88,580) -- (164,580) -- (164,650) -- (236,650) -- (306,650) -- (378,650)
    -- (378,580) -- (450,580) -- (450,510) -- (520,510) -- (520,440) -- (590,440)
    -- (590,368) -- (590,298) -- (590,224) -- (520,224) -- (520,155) -- (450,155)
    -- (378,155) -- (306,155) -- (306,224) -- (236,224) -- cycle;

  % ---------- disk outlines ----------
  \draw[line width=5.6] (390.8,334.5) circle (136.2);
  \draw[line width=5.6] (258.5,505) circle (120.5);

  % ---------- two plain slash strokes ----------
  \draw[line width=5.0] (53,118) -- (123,187);
  \draw[line width=5.0] (606,486) -- (656,620);

  % ---------- leader arrows ----------
  \draw[-{Triangle[length=20pt,width=15pt]},line width=3.2] (794,199) -- (532,281);
  \draw[-{Triangle[length=20pt,width=15pt]},line width=3.2] (802,288) -- (597,327);
  \draw[-{Triangle[length=20pt,width=15pt]},line width=3.2] (798,504) -- (631,537);
  \draw[-{Triangle[length=20pt,width=15pt]},line width=3.2] (805,601) -- (735,610);

  % ---------- pin canvas ----------
  \pgfresetboundingbox
  \path[use as bounding box] (0,0) rectangle (907,733);
\end{tikzpicture}
\end{document}
